\section{Typed fusion}
\renewcommand{\currentEdge}{fp-physics}

\begin{frame}{Functional Programming + Physics}
  \begin{center}
    \ResearchTriangle{fp-physics}
  \end{center}
\end{frame}

\begin{frame}{Little contact so far}
  Functional programming and programming-language theory
  have made \textbf{little contact} with computational physics.

  \vspace{0.8em}
  \begin{columns}[T]
    \begin{column}{0.48\textwidth}
      \textbf{Understandable}
      \begin{itemize}\setlength\itemsep{0.35em}
        \item Performance is the priority
        \item You want hardware control
        \item Imperative / low-level tools dominate
      \end{itemize}
    \end{column}
    \begin{column}{0.48\textwidth}
      \textbf{But appealing}
      \begin{itemize}\setlength\itemsep{0.35em}
        \item Specs of algorithms can be very rich
        \item Lots of structure to encode
        \item Types / abstraction can pull their weight
      \end{itemize}
    \end{column}
  \end{columns}
\end{frame}

\begin{frame}{A mathematically rich domain}
  \textbf{Example:} group representations
  (e.g.\ $SU(2)$ spins in quantum many-body physics).

  \vspace{0.7em}
  Symmetry-respecting maps (intertwiners) carry nontrivial structure:
  \begin{itemize}\setlength\itemsep{0.35em}
    \item Tensor products fuse: $\tfrac12\otimes\tfrac12 \;\cong\; 1\oplus 0$
    \item Data lives in \textbf{blocks} indexed by irreps
    \item Shapes are dictated by representation theory --- not by ad hoc bookkeeping
  \end{itemize}

  \vspace{0.6em}
  {\small Natural question: can the type system track this structure?}
\end{frame}

\begin{frame}[fragile]{Length-checked linear algebra}
  Start simple: dimensions in the types.
  \vspace{0.5em}
\begin{semiverbatim}
vectorExample :: C 3
vectorExample = vec (1,2,3)

tensorProductExample :: C 3 ⊗ C 2
tensorProductExample =
  vec (1,2,3) ⊗ vec (4,5)
\end{semiverbatim}
  \vspace{0.5em}
  Wrong length $\Rightarrow$ \textbf{compile-time} error --- not a silent
  reshape bug at runtime.
\end{frame}

\begin{frame}[fragile]{Irreps as types}
  Same idea, but the ``length'' is a representation:
  \vspace{0.5em}
\begin{semiverbatim}
irrepExample ::
  VectorOf SU2 (Irrep (Spin (1/2)))
irrepExample = vec (1,2)

directSumOfIrrepsExample ::
  Unfused SU2
    (Irrep (Spin (1/2))
       :⊕: Irrep (Spin (3/2)))
directSumOfIrrepsExample =
  (vec (1,2), vec (3,4,5,6))
\end{semiverbatim}
  \vspace{0.4em}
  {\small $\mathbb{C}^2$ for spin-$\tfrac12$; $\mathbb{C}^4$ for spin-$\tfrac32$ --- checked by the type.}
\end{frame}

\begin{frame}[fragile]{Fusion at the type level}
  Tensor products of irreps \textbf{fuse}:
  $\tfrac12\otimes\tfrac12 \;\cong\; 0\oplus 1$.
  \vspace{0.45em}
\begin{semiverbatim}
unfused :: Unfused SU2
  (Irrep (Spin (1/2)) :⊗: Irrep (Spin (1/2)))
unfused = mat ((1,2),(5,6))   -- ℂ² ⊗ ℂ²

fused :: Fused SU2
  (Irrep (Spin (1/2)) :⊗: Irrep (Spin (1/2)))
fused = (konst 1, vec (4,5,6))  -- ℂ¹ ⊕ ℂ³
\end{semiverbatim}
  \vspace{0.45em}
  Haskell computes the fused shape \textbf{at compile time};\\
  a wrong block size is rejected before you run anything.
\end{frame}

\begin{frame}[fragile]{Keeping only the trivial sector}
  Intertwiners live in the invariant ($j{=}0$) subspace ---
  Schur's lemma as a type:
  \vspace{0.45em}
\begin{semiverbatim}
trivialSectorExample ::
  Sym SU2
    (Irrep (Spin (1/2)) :⊗: Irrep (Spin (1/2)))
trivialSectorExample = konst 2   -- ℂ¹

-- 3/2 ⊗ 1/2 = 1 ⊕ 2  (no singlet)
trivialSectorExample2 ::
  Sym SU2
    (Irrep (Spin (3/2)) :⊗: Irrep (Spin (1/2)))
trivialSectorExample2 = EmptyFTrees
\end{semiverbatim}
\end{frame}

\begin{frame}[fragile]{From fusion to categories}
  Once objects are typed, morphisms compose by the
  \textbf{same} abstract law in every fusion category.
  \vspace{0.4em}
\begin{semiverbatim}
-- Mac Lane ladder on names ⌜f⌝, ⌜g⌝:
composeNames @Fib  nf ng
composeNames @Ising nf ng
\end{semiverbatim}
  \vspace{0.5em}
  \[
    \lceil f\rceil\otimes\lceil g\rceil
    \;\xrightarrow{\;\alpha,\;\varepsilon,\;\lambda\;}\;
    \lceil g\circ f\rceil
  \]
  \vspace{0.3em}
  Fibonacci and Ising share one implementation ---\\
  the mathematics supplies the abstraction boundary.
\end{frame}

\begin{frame}{Design principle}
  The mathematics shows how to \textbf{abstract} and
  \textbf{design} the software:
  \begin{itemize}\setlength\itemsep{0.45em}
    \item Representation theory $\;\Rightarrow\;$ block shapes in the types
    \item Fusion rules $\;\Rightarrow\;$ checked at compile time
    \item Compact closed structure $\;\Rightarrow\;$ one compositional API
      across theories (SU(2), Fib, Ising, \ldots)
  \end{itemize}

  \vspace{0.7em}
  {\small Code that looks like the maths --- not an imperative simulation
  with the same numeric answers.}
\end{frame}
